\chapter{Conclusion}

This project developed Verdant Search, a production-ready multimodal enterprise search engine that addresses critical gaps in information retrieval through integrated solutions spanning system architecture, ranking algorithms, and conversational AI. Deployed in collaboration with Infinity Datatech Sdn. Bhd., the system demonstrates that sophisticated multimodal reasoning can be achieved within the performance constraints of real-world enterprise environments.

\section{Project Summary}

The Verdant Search platform implements a complete end-to-end pipeline for enterprise-scale multimodal search, successfully addressing three core challenges identified in existing systems.

\textbf{Scalable Multimodal Architecture.} The system achieves sub-4-second P90 query latency on million-document collections through a two-stage ranking pipeline where fast hybrid retrieval (BM25 + HNSW vector search) identifies top-100 candidates, followed by expensive multimodal reranking applied only to the top-20 documents. Distributed crawling using Kafka coordinates document ingestion from heterogeneous enterprise sources (Confluence, Google Drive, Git, network shares), while dual indexing maintains both inverted indices and HNSW vector stores. The microservices architecture deployed on Kubernetes enables horizontal scaling and fault tolerance, supporting concurrent enterprise users in production deployment.

\textbf{Multimodal Listwise Reranking.} This work represents the first application of listwise learning-to-rank to genuinely multimodal documents. The BLIP-2-based reranker extends RankZephyr's setwise prompting to jointly evaluate textual and visual content, assessing cross-modal relationships including information consistency, contradiction detection, and complementarity. The sliding window approach reduces computational complexity from O(N²) pairwise comparisons to O(N) setwise evaluations while maintaining ranking quality, enabling effective ranking of business documents where charts, diagrams, and screenshots carry critical information.

\textbf{Iterative LLM-Guided Search.} The system implements a closed-loop architecture where search results drive LLM-generated summaries, which in turn refine subsequent retrieval through intelligent query rewriting. The dual-panel interface seamlessly integrates traditional document search with conversational AI, enabling users to iteratively refine information needs through natural dialogue. When users pose follow-up questions, the LLM rewrites queries incorporating conversation history, triggers fresh multimodal retrieval and reranking, and generates grounded answers with inline citations. This addresses the limitation of existing RAG systems that treat retrieval as one-shot preprocessing without feedback mechanisms.

\section{Research Gaps Addressed}

This work makes significant contributions to addressing critical gaps identified in the literature review:

\textbf{Gap 1: Multimodal Reranking.} All existing learning-to-rank methods from classical approaches to recent LLM-based rerankers operate exclusively on text. While MMDocIR demonstrated visual information value for first-stage retrieval, the reranking stage where cross-encoders prove most impactful remained unexplored for multimodal content. This project demonstrates that listwise reranking can be effectively extended to multimodal documents, with the vision-language reranker jointly evaluating textual and visual content to enable effective ranking of business documents where visual quality significantly impacts relevance.

\textbf{Gap 2: Multimodal RAG.} Current RAG systems assume text-only retrieval and generation, failing to leverage visual evidence that could improve answer quality for visually grounded questions. Verdant Search demonstrates bidirectional integration where multimodal retrieval improves AI-generated responses and LLM-generated insights refine subsequent search queries. The conversational interface enables users to ask about specific charts or diagrams, with the system retrieving and reasoning over both textual and visual content to generate grounded answers.

\textbf{Gap 3: Iterative Multimodal Search.} Existing search systems, even those with RAG capabilities, do not support iterative refinement where LLM reasoning improves retrieval quality. The closed-loop architecture implemented in this project enables progressive search refinement through conversational interaction, with each user question triggering query rewriting, fresh multimodal retrieval, reranking, and answer generation. This iterative approach better matches natural information-seeking behavior than one-shot retrieval.

\textbf{Gap 4: Enterprise Deployment and Scalability.} While academic research has explored multimodal retrieval in controlled settings, practical deployment at enterprise scale requires addressing engineering challenges including distributed coordination, incremental indexing, query latency optimization, and fault tolerance. This project demonstrates that sophisticated multimodal reasoning can be achieved within production performance constraints through strategic architectural choices: two-stage ranking to bound expensive inference, Redis caching for frequent queries, and horizontal scaling via stateless microservices.

\section{Practical Impact}

Deployment within Infinity Datatech addresses concrete operational challenges. Engineers previously spent 30-45 minutes daily searching across fragmented systems with an average of 7 context switches per search session. The unified interface reduces this to a single query point, with early user testing indicating a 60\% reduction in time-to-answer for common technical queries. The multimodal retrieval capability enables finding documents based on visual content descriptions, addressing the previous limitation where critical information in charts and diagrams was not searchable. The conversational AI interface allows support staff to refine queries through natural language during live customer interactions, reducing cognitive load while maintaining accuracy through inline citations.

\section{Limitations and Future Directions}

Current limitations include computational requirements (GPU acceleration limits concurrent users to approximately 50 on a single RTX 5090), optimization for English-language technical documentation (effectiveness on other languages or specialized domains not systematically evaluated), periodic batch index updates rather than real-time streaming, and evaluation reliance on MMDocIR which contains only 313 documents across 10 domains.

Future research directions include multimodal query understanding (accepting image inputs for visual similarity search), cross-lingual multimodal retrieval leveraging CLIP's language-agnostic visual representations, personalized ranking incorporating user-specific signals while preserving privacy, explainable multimodal ranking highlighting which textual passages and visual elements contributed to decisions, efficient model compression through distillation and quantization to improve scalability, and development of large-scale multimodal benchmarks with comprehensive relevance judgments to enable rigorous evaluation.

\section{Final Remarks}

This project demonstrates that multimodal search systems combining sophisticated vision-language reasoning with conversational AI capabilities can be successfully deployed in enterprise environments. The key insight is that effective multimodal search requires integrated solutions spanning multiple system layers: distributed infrastructure for scalable indexing, hybrid retrieval combining lexical and semantic methods, listwise reranking with cross-modal reasoning, and conversational interfaces enabling iterative refinement. The collaboration with Infinity Datatech exemplifies how academic research can address industry needs while advancing the state of the art, delivering immediate practical value while opening new research directions for the information retrieval community.
